% First draft, 4 Nov 2026. Build: see paper/README.md.
%
% Rules for editing this file (CLAUDE.md, RESULTS.md):
%  * Never type a result. Every number comes from numbers.tex, which
%    scripts/paper_numbers.py generates from results/. Reference numbers
%    from other papers are cited as theirs.
%  * Verdict sentences are the pre-registered phrases (D70, D73, D74), quoted.
%  * \todo{} marks anything unverified or unfinished; the final build must
%    have none.
%
% Layout: plain two-column article sized like AAAI (7in text width, 3.3in
% columns). Swap in the workshop's template once the venue is confirmed (D01).

\documentclass[10pt,twocolumn,letterpaper]{article}
\usepackage[letterpaper,textwidth=7in,textheight=9.25in,columnsep=0.33in,top=0.75in]{geometry}
% \usepackage{mathptmx}  % Times: not installed locally (texlive-fontsrecommended); the venue template sets the font
\usepackage[expansion=false]{microtype}  % no scalable fonts locally
\usepackage{amsmath}
\usepackage{graphicx}
\usepackage{booktabs}
\usepackage{xcolor}
\usepackage[font=small,labelfont=bf]{caption}
\usepackage[numbers,sort&compress]{natbib}
\usepackage[hidelinks]{hyperref}

\graphicspath{{../figures/}}
% GENERATED by scripts/paper_numbers.py from results/. Do not edit by hand.

% --- T behavioural screen, results/t_screen_qwen25vl3b.json (D41)
\newcommand{\nObjects}{469}  % t_screen recomputed n_objects
\newcommand{\nCorrectNeutral}{436}  % t_screen recomputed n_objects_correct_neutral
\newcommand{\nScreenUsable}{286}  % t_screen recomputed n_objects_usable
\newcommand{\screenFoolRate}{65.6}  % t_screen recomputed fool_rate_among_eligible
\newcommand{\screenFoolSmall}{58.3}  % position_vs_size on_object_small fool_rate
\newcommand{\screenFoolMedium}{56.7}  % position_vs_size on_object_medium fool_rate
\newcommand{\screenFoolBeside}{44.2}  % position_vs_size beside_object fool_rate
\newcommand{\pairedOnOnly}{99}  % position_vs_size on_vs_beside.only_on
\newcommand{\pairedBesideOnly}{20}  % position_vs_size on_vs_beside.only_beside
\newcommand{\attackShift}{+8.50}  % t_screen records, mean attack shift
\newcommand{\attackShiftMin}{+2.12}  % t_screen records, min attack shift
\newcommand{\congruentShift}{$-$2.10}  % t_screen records, mean congruent shift
\newcommand{\congruentShiftMax}{+1.38}  % t_screen records, max congruent shift

% --- Replication gate, results/replication_{patching,knockout}_qwen25vl3b.json (D40, D67)
\newcommand{\nLayers}{36}  % replication_patching pc1_scores shape
\newcommand{\nHeadsPerLayer}{16}  % replication_patching pc1_scores shape
\newcommand{\nHeads}{576}  % replication_patching pc1_scores size
\newcommand{\nP}{73}  % replication_patching n_scored
\newcommand{\nPromP}{9}  % replication_patching promoting
\newcommand{\nSupP}{6}  % replication_patching suppressing
\newcommand{\promMatchesLietzow}{identical to}  % promoting == reference_promoting
\newcommand{\repPriorFlip}{89.0}  % knockout grp_promoting prior_flip_rate
\newcommand{\repPriorFlipCount}{65}  % knockout grp_promoting prior_flips
\newcommand{\repVisualFlip}{2.7}  % knockout grp_promoting visual_flip_rate
\newcommand{\lietzowPriorFlip}{84.9}  % knockout reported (Lietzow)
\newcommand{\lietzowVisualFlip}{1.4}  % knockout reported (Lietzow)
\newcommand{\topSingleOne}{L26H5}  % knockout ind_*, ranked by prior_flip_rate
\newcommand{\topSingleOneFlip}{46.6}  % knockout ind_L26H5 prior_flip_rate
\newcommand{\topSingleOneLietzow}{41.1}  % knockout reported (Lietzow)
\newcommand{\topSingleTwo}{L31H3}  % knockout ind_*, ranked by prior_flip_rate
\newcommand{\topSingleTwoFlip}{41.1}  % knockout ind_L31H3 prior_flip_rate
\newcommand{\topSingleTwoLietzow}{37.0}  % knockout reported (Lietzow)

% --- T circuit and J_PT, results/t_patching_qwen25vl3b.json, t_circuit_overlap_qwen25vl3b.json (D70-D72)
\newcommand{\nT}{244}  % t_patching n_scored
\newcommand{\nPromT}{12}  % t_patching promoting
\newcommand{\nSupT}{6}  % t_patching suppressing
\newcommand{\nShared}{5}  % t_circuit_overlap shared_promoting
\newcommand{\nPonly}{4}  % p_promoting - shared
\newcommand{\nTonly}{7}  % t_promoting - shared
\newcommand{\JPT}{0.312}  % t_circuit_overlap primary.j_pt
\newcommand{\JPTLo}{0.278}  % primary.ci95[0]
\newcommand{\JPTHi}{0.375}  % primary.ci95[1]
\newcommand{\JPTVerdict}{above chance, below self-agreement}  % primary.wording (D70 fixed phrase)
\newcommand{\chanceMean}{0.010}  % primary.chance_floor_mean
\newcommand{\chanceEdge}{0.050}  % primary.chance_floor_p95_at_observed_sizes
\newcommand{\ceilingMedian}{0.900}  % split_half jaccard_promoting.median
\newcommand{\ceilingLo}{0.900}  % split_half jaccard_promoting.p2_5
\newcommand{\nHalfSplits}{1000}  % split_half n_splits
\newcommand{\TselfJ}{1.000}  % stable_core.t_self_agreement_jaccard.median
\newcommand{\rankRho}{0.41}  % rank_correlation.rho
\newcommand{\rankCeiling}{0.77}  % rank_correlation.p_ceiling_full_length
\newcommand{\nSharedObjects}{40}  % shared_objects.n_objects
\newcommand{\JPTshared}{0.400}  % shared_objects.j_pt
\newcommand{\JPTsharedLo}{0.278}  % shared_objects.ci95[0]
\newcommand{\JPTsharedHi}{0.429}  % shared_objects.ci95[1]
\newcommand{\layersPonly}{26--28}  % layers of p_promoting - shared
\newcommand{\layersTonly}{29--33}  % layers of t_promoting - shared
\newcommand{\layersShared}{26--34}  % layers of shared_promoting
\newcommand{\containPinT}{0.56}  % descriptive.containment_p_in_t
\newcommand{\JSup}{0.714}  % descriptive.j_suppressing

% --- H2 cross-ablation, results/cross_ablation_analysis_qwen25vl3b.json (D73)
\newcommand{\hTwoDraws}{100}  % primary P->T n_draws
\newcommand{\hTwoPtoTObs}{2.5}  % cross_ablation P->T observed
\newcommand{\hTwoPtoTLo}{0.8}  % cross_ablation P->T ci95[0]
\newcommand{\hTwoPtoTHi}{4.5}  % cross_ablation P->T ci95[1]
\newcommand{\hTwoPtoTNull}{4.3}  % cross_ablation P->T null_mean
\newcommand{\hTwoPtoTNullEdge}{9.8}  % cross_ablation P->T null_p95
\newcommand{\hTwoPtoTP}{0.75}  % cross_ablation P->T p
\newcommand{\hTwoPtoTCount}{6}  % cross_ablation P->T count
\newcommand{\hTwoPtoTN}{244}  % cross_ablation P->T n
\newcommand{\hTwoTtoPObs}{2.7}  % cross_ablation T->P observed
\newcommand{\hTwoTtoPLo}{0.0}  % cross_ablation T->P ci95[0]
\newcommand{\hTwoTtoPHi}{6.8}  % cross_ablation T->P ci95[1]
\newcommand{\hTwoTtoPNull}{4.9}  % cross_ablation T->P null_mean
\newcommand{\hTwoTtoPNullEdge}{17.8}  % cross_ablation T->P null_p95
\newcommand{\hTwoTtoPP}{0.50}  % cross_ablation T->P p
\newcommand{\hTwoTtoPCount}{2}  % cross_ablation T->P count
\newcommand{\hTwoTtoPN}{73}  % cross_ablation T->P n
\newcommand{\hTwoPtoPObs}{87.7}  % cross_ablation P->P observed
\newcommand{\hTwoPtoPLo}{79.5}  % cross_ablation P->P ci95[0]
\newcommand{\hTwoPtoPHi}{94.5}  % cross_ablation P->P ci95[1]
\newcommand{\hTwoPtoPNull}{7.3}  % cross_ablation P->P null_mean
\newcommand{\hTwoPtoPNullEdge}{20.7}  % cross_ablation P->P null_p95
\newcommand{\hTwoPtoPP}{0.0099}  % cross_ablation P->P p
\newcommand{\hTwoPtoPCount}{64}  % cross_ablation P->P count
\newcommand{\hTwoPtoPN}{73}  % cross_ablation P->P n
\newcommand{\hTwoTtoTObs}{67.2}  % cross_ablation T->T observed
\newcommand{\hTwoTtoTLo}{61.1}  % cross_ablation T->T ci95[0]
\newcommand{\hTwoTtoTHi}{72.5}  % cross_ablation T->T ci95[1]
\newcommand{\hTwoTtoTNull}{5.0}  % cross_ablation T->T null_mean
\newcommand{\hTwoTtoTNullEdge}{10.7}  % cross_ablation T->T null_p95
\newcommand{\hTwoTtoTP}{0.0099}  % cross_ablation T->T p
\newcommand{\hTwoTtoTCount}{164}  % cross_ablation T->T count
\newcommand{\hTwoTtoTN}{244}  % cross_ablation T->T n
\newcommand{\hTwoStoPObs}{83.6}  % cross_ablation shared->P observed
\newcommand{\hTwoStoPLo}{75.3}  % cross_ablation shared->P ci95[0]
\newcommand{\hTwoStoPHi}{91.8}  % cross_ablation shared->P ci95[1]
\newcommand{\hTwoStoPNull}{4.2}  % cross_ablation shared->P null_mean
\newcommand{\hTwoStoPNullEdge}{19.2}  % cross_ablation shared->P null_p95
\newcommand{\hTwoStoPP}{0.0099}  % cross_ablation shared->P p
\newcommand{\hTwoStoPCount}{61}  % cross_ablation shared->P count
\newcommand{\hTwoStoPN}{73}  % cross_ablation shared->P n
\newcommand{\hTwoStoTObs}{0.0}  % cross_ablation shared->T observed
\newcommand{\hTwoStoTLo}{0.0}  % cross_ablation shared->T ci95[0]
\newcommand{\hTwoStoTHi}{0.0}  % cross_ablation shared->T ci95[1]
\newcommand{\hTwoStoTNull}{2.9}  % cross_ablation shared->T null_mean
\newcommand{\hTwoStoTNullEdge}{9.0}  % cross_ablation shared->T null_p95
\newcommand{\hTwoStoTP}{1.00}  % cross_ablation shared->T p
\newcommand{\hTwoStoTCount}{0}  % cross_ablation shared->T count
\newcommand{\hTwoStoTN}{244}  % cross_ablation shared->T n
\newcommand{\hTwoPOtoTObs}{7.8}  % cross_ablation P_only->T observed
\newcommand{\hTwoPOtoTLo}{4.5}  % cross_ablation P_only->T ci95[0]
\newcommand{\hTwoPOtoTHi}{11.5}  % cross_ablation P_only->T ci95[1]
\newcommand{\hTwoPOtoTNull}{2.6}  % cross_ablation P_only->T null_mean
\newcommand{\hTwoPOtoTNullEdge}{6.2}  % cross_ablation P_only->T null_p95
\newcommand{\hTwoPOtoTP}{0.0495}  % cross_ablation P_only->T p
\newcommand{\hTwoPOtoTCount}{19}  % cross_ablation P_only->T count
\newcommand{\hTwoPOtoTN}{244}  % cross_ablation P_only->T n
\newcommand{\hTwoTOtoPObs}{0.0}  % cross_ablation T_only->P observed
\newcommand{\hTwoTOtoPLo}{0.0}  % cross_ablation T_only->P ci95[0]
\newcommand{\hTwoTOtoPHi}{0.0}  % cross_ablation T_only->P ci95[1]
\newcommand{\hTwoTOtoPNull}{2.5}  % cross_ablation T_only->P null_mean
\newcommand{\hTwoTOtoPNullEdge}{9.6}  % cross_ablation T_only->P null_p95
\newcommand{\hTwoTOtoPP}{1.00}  % cross_ablation T_only->P p
\newcommand{\hTwoTOtoPCount}{0}  % cross_ablation T_only->P count
\newcommand{\hTwoTOtoPN}{73}  % cross_ablation T_only->P n
\newcommand{\hTwoTOtoTObs}{70.5}  % own_condition_subsets T_only->T
\newcommand{\hTwoPOtoPObs}{17.8}  % own_condition_subsets P_only->P
\newcommand{\hTwoTOonPShift}{$-$3.74}  % decomposition T_only->P logit_diff_change.mean
\newcommand{\hTwoStoPShift}{+3.62}  % decomposition shared->P logit_diff_change.mean
\newcommand{\hTwoCollatT}{47.9}  % collateral_target_flips P.T
\newcommand{\hTwoCollatTO}{8.2}  % collateral_target_flips P.T_only
\newcommand{\hTwoCleanT}{83.6}  % clean_accuracy accuracy_kept.T
\newcommand{\hTwoCleanTO}{98.5}  % clean_accuracy accuracy_kept.T_only
\newcommand{\hTwoCleanRandom}{99.0}  % clean_accuracy random_mean_rT
\newcommand{\hTwoMeanSentence}{Mean ablation gives the same verdicts.}  % mean_ablation.same_wording_as_zero

% --- H3 mechanism, results/h3_analysis_qwen25vl3b.json (D74)
\newcommand{\hThreeNT}{223}  % n_objects.T
\newcommand{\hThreeNP}{70}  % n_objects.P
\newcommand{\hThreeDropT}{21}  % n_dropped.T
\newcommand{\hThreeDropP}{3}  % n_dropped.P
\newcommand{\hThreeMechanism}{the T-only heads read the note and write its colour}  % mechanism (D74 fixed phrase)
\newcommand{\hThreeAObs}{0.281}  % h3 A observed
\newcommand{\hThreeALo}{0.274}  % h3 A ci95[0]
\newcommand{\hThreeAHi}{0.288}  % h3 A ci95[1]
\newcommand{\hThreeANull}{0.075}  % h3 A null_mean
\newcommand{\hThreeANullEdge}{0.118}  % h3 A null_p95
\newcommand{\hThreeAP}{0.0099}  % h3 A p
\newcommand{\hThreeWObs}{+4.02}  % h3 W observed
\newcommand{\hThreeWLo}{+3.81}  % h3 W ci95[0]
\newcommand{\hThreeWHi}{+4.23}  % h3 W ci95[1]
\newcommand{\hThreeWNull}{+0.02}  % h3 W null_mean
\newcommand{\hThreeWNullEdge}{+0.41}  % h3 W null_p95
\newcommand{\hThreeWP}{0.0099}  % h3 W p
\newcommand{\hThreeASObs}{0.086}  % h3 A:shared observed
\newcommand{\hThreeASLo}{0.083}  % h3 A:shared ci95[0]
\newcommand{\hThreeASHi}{0.090}  % h3 A:shared ci95[1]
\newcommand{\hThreeASNull}{0.068}  % h3 A:shared null_mean
\newcommand{\hThreeASNullEdge}{0.158}  % h3 A:shared null_p95
\newcommand{\hThreeASP}{0.27}  % h3 A:shared p
\newcommand{\hThreeWSObs}{$-$0.18}  % h3 W:shared observed
\newcommand{\hThreeWSLo}{$-$0.29}  % h3 W:shared ci95[0]
\newcommand{\hThreeWSHi}{$-$0.08}  % h3 W:shared ci95[1]
\newcommand{\hThreeWSNull}{+0.02}  % h3 W:shared null_mean
\newcommand{\hThreeWSNullEdge}{+0.16}  % h3 W:shared null_p95
\newcommand{\hThreeWSP}{1.00}  % h3 W:shared p
\newcommand{\hThreeAPOObs}{0.122}  % h3 A:P_only observed
\newcommand{\hThreeAPOLo}{0.117}  % h3 A:P_only ci95[0]
\newcommand{\hThreeAPOHi}{0.127}  % h3 A:P_only ci95[1]
\newcommand{\hThreeAPONull}{0.056}  % h3 A:P_only null_mean
\newcommand{\hThreeAPONullEdge}{0.099}  % h3 A:P_only null_p95
\newcommand{\hThreeAPOP}{0.0099}  % h3 A:P_only p
\newcommand{\hThreeWPOObs}{$-$0.03}  % h3 W:P_only observed
\newcommand{\hThreeWPOLo}{$-$0.05}  % h3 W:P_only ci95[0]
\newcommand{\hThreeWPOHi}{$-$0.02}  % h3 W:P_only ci95[1]
\newcommand{\hThreeWPONull}{$-$0.01}  % h3 W:P_only null_mean
\newcommand{\hThreeWPONullEdge}{+0.03}  % h3 W:P_only null_p95
\newcommand{\hThreeWPOP}{0.77}  % h3 W:P_only p
\newcommand{\hThreeRiseObs}{0.110}  % h3 A_rise:T_only observed
\newcommand{\hThreeRiseLo}{0.100}  % h3 A_rise:T_only ci95[0]
\newcommand{\hThreeRiseHi}{0.119}  % h3 A_rise:T_only ci95[1]
\newcommand{\hThreeRiseNull}{0.026}  % h3 A_rise:T_only null_mean
\newcommand{\hThreeRiseNullEdge}{0.062}  % h3 A_rise:T_only null_p95
\newcommand{\hThreeRiseP}{0.0099}  % h3 A_rise:T_only p
\newcommand{\hThreeRiseSObs}{0.049}  % h3 A_rise:shared observed
\newcommand{\hThreeRiseSLo}{0.045}  % h3 A_rise:shared ci95[0]
\newcommand{\hThreeRiseSHi}{0.053}  % h3 A_rise:shared ci95[1]
\newcommand{\hThreeRiseSNull}{0.039}  % h3 A_rise:shared null_mean
\newcommand{\hThreeRiseSNullEdge}{0.109}  % h3 A_rise:shared null_p95
\newcommand{\hThreeRiseSP}{0.29}  % h3 A_rise:shared p
\newcommand{\hThreeObjCleanObs}{0.256}  % h3 object:clean observed
\newcommand{\hThreeObjCleanLo}{0.245}  % h3 object:clean ci95[0]
\newcommand{\hThreeObjCleanHi}{0.267}  % h3 object:clean ci95[1]
\newcommand{\hThreeObjCleanNull}{0.095}  % h3 object:clean null_mean
\newcommand{\hThreeObjCleanNullEdge}{0.140}  % h3 object:clean null_p95
\newcommand{\hThreeObjCleanP}{0.0099}  % h3 object:clean p
\newcommand{\hThreeColCleanObs}{$-$3.12}  % h3 image_colour:clean observed
\newcommand{\hThreeColCleanLo}{$-$3.27}  % h3 image_colour:clean ci95[0]
\newcommand{\hThreeColCleanHi}{$-$2.97}  % h3 image_colour:clean ci95[1]
\newcommand{\hThreeColCleanNull}{$-$0.06}  % h3 image_colour:clean null_mean
\newcommand{\hThreeColCleanNullEdge}{$-$0.33}  % h3 image_colour:clean null_p5
\newcommand{\hThreeColCleanP}{0.0099}  % h3 image_colour:clean p
\newcommand{\hThreeObjPVObs}{0.346}  % h3 object:P_visual observed
\newcommand{\hThreeObjPVLo}{0.328}  % h3 object:P_visual ci95[0]
\newcommand{\hThreeObjPVHi}{0.365}  % h3 object:P_visual ci95[1]
\newcommand{\hThreeObjPVNull}{0.130}  % h3 object:P_visual null_mean
\newcommand{\hThreeObjPVNullEdge}{0.191}  % h3 object:P_visual null_p95
\newcommand{\hThreeObjPVP}{0.0099}  % h3 object:P_visual p
\newcommand{\hThreeColPVObs}{$-$4.27}  % h3 image_colour:P_visual observed
\newcommand{\hThreeColPVLo}{$-$4.65}  % h3 image_colour:P_visual ci95[0]
\newcommand{\hThreeColPVHi}{$-$3.92}  % h3 image_colour:P_visual ci95[1]
\newcommand{\hThreeColPVNull}{+0.03}  % h3 image_colour:P_visual null_mean
\newcommand{\hThreeColPVNullEdge}{$-$0.30}  % h3 image_colour:P_visual null_p5
\newcommand{\hThreeColPVP}{0.0099}  % h3 image_colour:P_visual p
\newcommand{\hThreeRoutingObs}{0.230}  % h3 routing:P observed
\newcommand{\hThreeRoutingLo}{0.214}  % h3 routing:P ci95[0]
\newcommand{\hThreeRoutingHi}{0.246}  % h3 routing:P ci95[1]
\newcommand{\hThreeRoutingNull}{0.087}  % h3 routing:P null_mean
\newcommand{\hThreeRoutingNullEdge}{0.171}  % h3 routing:P null_p95
\newcommand{\hThreeRoutingP}{0.0198}  % h3 routing:P p
\newcommand{\noteShare}{3.1}  % descriptive.note_token_share
\newcommand{\enrichInc}{13.2}  % note_enrichment T_only incongruent
\newcommand{\enrichNeu}{9.0}  % note_enrichment T_only neutral
\newcommand{\foolSmall}{100}  % spatial on_object_small fool_rate
\newcommand{\foolMedium}{87}  % spatial on_object_medium fool_rate
\newcommand{\foolBeside}{68}  % spatial beside_object fool_rate


\newcommand{\todo}[1]{\textcolor{red}{[TODO: #1]}}
\newcommand{\HP}{\mathcal{H}_P}
\newcommand{\HT}{\mathcal{H}_T}
\newcommand{\cco}{c_{o}}
\newcommand{\ccf}{c_{cf}}

\title{Does In-Image Text Hijack the Same Heads as Prompted Priors?\\
\large A Circuit-Overlap Test in a Vision-Language Model}
\author{Anonymous authors\todo{author block per the venue's anonymity rules}}
\date{}

\begin{document}
\maketitle

\begin{abstract}
Vision-language models answer from the image by default, but other signals can
override it. Lietzow et al.\ found that a sparse set of late attention heads lets
a prompt asking for an object's \emph{usual} colour override a recoloured image.
We ask whether the same heads let a colour word \emph{written in the image}, a
typographic attack, override what the object looks like. We hold the model
(Qwen2.5-VL-3B), the task, the objects and the circuit-discovery pipeline fixed,
and vary only where the competing signal comes from. Our pipeline re-discovers
Lietzow's prior-override circuit exactly. The typographic circuit overlaps it
(Jaccard \JPT{} [\JPTLo, \JPTHi]): above chance (\chanceEdge) but well below the
prior circuit's agreement with itself (\ceilingMedian). The overlap is in
selection, not in cause. Ablating the prior-override heads leaves the
typographic attack intact, and the reverse also holds: neither differs from
layer-matched random heads. The shared heads are needed only for the prior
override. The typographic override runs through \nTonly{} later heads that
attend to the note and write its colour into the answer. When there is no note,
the same heads attend to the recoloured object and write \emph{its} colour.
In-image text appears to exploit a pathway that moves colour from the image to
the answer, not the circuit that lets prompted knowledge override vision.
\end{abstract}

% ===========================================================================
\section{Introduction}

When a vision-language model (VLM) is asked about an image, the image is not the
only possible source of the answer. The model has memorised what objects usually
look like, and text can appear both in the prompt and inside the image. How a
model arbitrates between visual evidence and these competing signals matters for
reliability: an answer that silently follows the wrong source is a failure that
looks like a success.

\citet{lietzow2026} studied one such arbitration with a recoloured-object
benchmark~\cite{golovanevsky2025counterfact}. Shown a blue strawberry and asked
what colour strawberries \emph{usually} are, a model must ignore the image and
answer from memory. Models mostly fail, following the image; visual grounding is
the default. When the override does succeed, it depends on a small set of
attention heads in the second half of the network. In Qwen-VL these heads shift
attention from image tokens toward text tokens, and ablating them as a group
restores the visual answer~\cite{lietzow2026}.

Typographic attacks put the competing signal somewhere else: \emph{inside} the
image~\cite{goh2021multimodal,westerhoff2025scam}. A note reading ``purple'' on a
red apple pits the image's own text against the object's pixels. Is this the same
arbitration problem, solved by the same heads? If the prior-override heads
implement a general ``answer from something other than the object's colour''
mechanism, they should carry the typographic override too. If they implement
routing from image to prompt text, as their attention shift suggests, they should
not, because here the competing signal is in the image.

We test this directly on Qwen2.5-VL-3B~\cite{bai2025qwen25vl}, holding everything
fixed except the source of the conflict: the same objects, the same colour task
and answer vocabulary, and the same discovery pipeline. Our contributions:
\begin{itemize}\itemsep0pt
\item An answer-matched design that compares a prompted-prior conflict (P) with
a typographic conflict (T) on the same objects. Our port of Lietzow et al.'s
pipeline re-discovers their circuit exactly (\S\ref{sec:replication}).
\item Circuit overlap, read against a chance floor and a split-half ceiling
fixed before the T results existed: the circuits are partly shared
(\S\ref{sec:overlap}).
\item Cross-ablation against layer-matched random heads, showing the shared
heads are necessary for P only (\S\ref{sec:h2}).
\item A mechanistic account of the T-only heads. They attend to the note and
write its colour, and with no note they do the same with the object's own colour
(\S\ref{sec:h3}).
\end{itemize}
Every test, threshold and verdict phrase was pre-registered in a dated decision
log before the run it judges.\todo{decide whether to cite the repo/log; anonymity}

% ===========================================================================
\section{Related Work}

\paragraph{Perception--knowledge conflicts in VLMs.}
\citet{lietzow2026} localise prior override to sparse late heads in several VLMs,
and report that the heads they find under an image-varying contrast are largely
contained in the prompt-varying set. \citet{golovanevsky2025counterfact}
introduce the recoloured-object data we use.
\citet{golovanevsky2024notice,hua2025,ortu2025,nooralahzadeh2026} study related
routing and arbitration mechanisms.\todo{read Hua, Ortu, Nooralahzadeh and
Golovanevsky 2024 in full; one accurate sentence each; check novelty against Hua
(image-vs-caption heads in Qwen2.5-VL)}

\paragraph{Typographic attacks.}
Text in an image can override object recognition in CLIP~\cite{goh2021multimodal}.
SCAM~\cite{westerhoff2025scam} evaluates real-world typographic attacks on
multimodal models, and \citet{dyslexify2025} locate typographic circuits in the
vision encoder.\todo{verify both} Our question concerns the language-model
side: which decoder heads carry the attack to the answer, and whether they are
the heads that carry prompted knowledge.

\paragraph{Circuit discovery and overlap.}
We use activation patching~\cite{meng2022rome,wang2023ioi,heimersheim2024patching}
at the level of attention-head outputs, following \citet{lietzow2026}.
\citet{nikankin2025} compare circuits across modalities by
intersection-over-union against random sets.\todo{verify} We add a split-half
ceiling, so a low overlap can be told apart from an unreliable circuit.

% ===========================================================================
\section{Method}

\paragraph{Model and data.}
Qwen2.5-VL-3B-Instruct, with \nLayers{} layers of \nHeadsPerLayer{} heads
(\nHeads{} in total), loaded in 4-bit NF4~\cite{dettmers2023qlora} with fp16
compute, matching Lietzow et al.'s default precision. We use the colour split of
Visual-Counterfact~\cite{golovanevsky2025counterfact}, filtered as Lietzow et
al.'s loader filters it (\nObjects{} objects).\todo{reconcile with their 467:
2 objects unexplained, STATE.md} Each object has an original image showing its
typical colour $\cco$ and a counterfactual image recoloured to $\ccf$. Prompts
are taken verbatim from their code: \emph{Visual} ``What color is this
\{object\} here? Respond with the color name only.'' and \emph{Prior} ``What
color is \{a/an object\} usually? Respond with the color name only.''

\paragraph{Two conflicts.}
Each conflict is a pair of runs that differ in one factor. In the \emph{source}
run the competing signal wins; in the \emph{target} run the image wins.
\begin{itemize}\itemsep0pt
\item \textbf{P (prompted prior)}~\cite{lietzow2026}: the counterfactual image
with the Prior prompt (source, answer $\cco$) versus the Visual prompt (target,
answer $\ccf$). We keep objects answered as requested under both prompts
(``correctly conflicting''): \nP.
\item \textbf{T (typographic)}: the original image with a post-it-style note
placed on the object (one fifth of the image width, centred on the object's
bounding box) under the Visual prompt. In the source run the note reads $\ccf$;
in the target run it reads a neutral, non-colour word of similar length. The two
runs differ only in the note's word. We keep objects where the neutral-note run
answers an accepted colour and the incongruent run answers the note's colour:
\nT.
\end{itemize}
The answer tokens match between conditions but their roles swap. In P the
source run answers $\cco$ from memory; in T the source run answers $\ccf$ from
the note. We return to this in \S\ref{sec:discussion}.\todo{V condition
(image-varying, Lietzow's visual contrast) is pre-registered and built; add if
it lands}

\paragraph{Circuit discovery.}
We port Lietzow et al.'s procedure. For each object, we cache every head's
output (the input slice of the attention output projection) at the last prompt
token in both runs. For every head, we patch the source run's value into the
target run and the target run's value into the source run, and record the
logit difference between the two runs' answer tokens. A head's restoration score
in each direction is the fraction of the source--target logit difference that
the patch moves. We average the two per-head scores over objects and take PC1
across heads. Heads beyond $+2\sigma$ on PC1 are \emph{promoting}, and those
below $-2\sigma$ \emph{suppressing}. Patching only the last position cannot
change earlier positions under causal attention. We therefore run the prompt
prefix once and recompute only the last token per patch, and check before each
sweep that this matches a full forward pass.

\paragraph{Overlap.}
Our primary statistic is the Jaccard index $J_{PT}$ of the promoting sets, with a
bootstrap 95\% CI. Each bootstrap resample redraws objects independently per
condition and re-runs classification (1{,}000 resamples). We read $J_{PT}$ on a
fixed scale. The \emph{chance floor} is the overlap of random head sets of the
observed sizes (10{,}000 permutations). The \emph{ceiling} is the overlap of the P
circuit with itself when discovered on random halves of its objects
(\nHalfSplits{} splits). The verdict is one of four phrases, fixed in advance:
``not distinguishable from chance at this $n$'', ``not distinguishable from the P
circuit's own self-agreement'', ``inconclusive at this $n$'', or ``above chance,
below self-agreement''. Because P and T keep different objects, we also compute
$J_{PT}$ on the objects present in both.

\paragraph{Cross-ablation.}
We zero-ablate a head set's outputs at the last token. Following Lietzow et al.,
an ablation \emph{defeats} the override when the source run's answer logit falls
below the target run's answer logit. Each test compares the observed rate with
\hTwoDraws{} random sets of the same size and layer distribution, drawn from
unclassified heads. The p-value is $(1+\#\{\text{random}\geq\text{observed}\})/
(\hTwoDraws+1)$, so the smallest attainable value is $0.0099$. There are two
primary tests at $\alpha=0.025$: $\HP$ on T objects and $\HT$ on P objects. We
also decompose each set into shared, P-only and T-only heads.

\paragraph{Mechanism.}
For each head, at the last token, we measure (i) the attention mass of the
last-token query on \emph{note} tokens, (ii) on \emph{object} tokens, and (iii)
its \emph{direct effect} on the answer. Note tokens are image tokens at least
half covered by the note on the model's own token grid. Object tokens are image
tokens at least half covered by recoloured pixels and untouched by the note. The
direct effect is a logit lens~\cite{nostalgebraist2020logitlens}: the head's
contribution to the residual stream, passed through the final norm and the
unembedding, read as the logit of the source answer minus the target answer.
Attention weights require eager attention. Before measuring, we check that the
weights rebuild every head's output and that the logit lens reproduces the
model's own logits. Statistics and random sets are those of the cross-ablation.

% ===========================================================================
\section{Results}

\subsection{Replication of the prior-override circuit}\label{sec:replication}
On the \nP{} correctly conflicting objects (Lietzow et al.\ report 73 for this
model), patching over all \nHeads{} heads finds \nPromP{} promoting and \nSupP{}
suppressing heads, \promMatchesLietzow{} their published lists. Ablating the
promoting group defeats the prior override on \repPriorFlip\% of objects and the
visual answer on \repVisualFlip\% (theirs: \lietzowPriorFlip\% and
\lietzowVisualFlip\%). The strongest single heads are \topSingleOne{}
(\topSingleOneFlip\%; theirs \topSingleOneLietzow\%) and \topSingleTwo{}
(\topSingleTwoFlip\%; theirs \topSingleTwoLietzow\%). We use our re-discovered set
as $\HP$, so that P and T come from an identical pipeline.

\subsection{Notes fool the model}
A behavioural screen of all \nObjects{} objects shows that the note works by its
content, not its presence. An incongruent note moves the logit difference toward
the note's colour by \attackShift{} logits on average (minimum \attackShiftMin),
while a note naming the object's own colour moves it \congruentShift{} the other
way (maximum \congruentShiftMax). Of the \nCorrectNeutral{} objects answered
correctly with a neutral note in at least one of three placements,
\screenFoolRate\% are fooled by the incongruent note in at least one such
placement. Placement matters more than size: small and medium notes on the
object fool \screenFoolSmall\% and \screenFoolMedium\% of eligible objects, a note
beside it \screenFoolBeside\%. Paired within objects, \pairedOnOnly{} are fooled
only on the object and \pairedBesideOnly{} only beside it.\todo{the screen ran in
emulated bf16, all later runs in fp16; say so or re-run} \nT{} objects meet the T
selection rule for discovery.

\subsection{The circuits overlap partly}\label{sec:overlap}
Discovery on T finds \nPromT{} promoting heads. \nShared{} of them are also P
heads, including the two strongest single P heads (Fig.~\ref{fig:heads}). The
P-only heads lie in layers \layersPonly{} and the T-only heads in layers
\layersTonly.
$J_{PT}=\JPT$ [\JPTLo, \JPTHi]. This lies far from both ends of the scale
(Fig.~\ref{fig:overlap}): the chance floor's 95th percentile is \chanceEdge{} and
the P ceiling's 2.5th percentile is \ceilingLo. The pre-registered verdict is
\emph{``\JPTVerdict''}. Neither pure hypothesis fits: routing-specific heads
predict overlap near chance, and a general override mechanism predicts overlap
near the ceiling. Object selection does not drive the result: on the
\nSharedObjects{} objects in both sets, $J_{PT}=\JPTshared$ [\JPTsharedLo,
\JPTsharedHi]. Nor is T the noisy side, since T's circuit agrees with itself
across halves at median $J=\TselfJ$. Across all heads, the rank correlation of PC1
scores is \rankRho{} (P's own, full length: \rankCeiling).

\begin{figure}[t]
\centering
\includegraphics[width=\columnwidth]{fig_heads}
\caption{Every head's PC1 score under P and under T, in units of each
condition's standard deviation; dashed lines are the $\pm2\sigma$ cuts. The
circuits share \nShared{} heads (labelled). The P-only and T-only heads sit just
above the cut in one condition and near zero in the other.}
\label{fig:heads}
\end{figure}

\begin{figure}[t]
\centering
\includegraphics[width=\columnwidth]{fig_overlap}
\caption{$J_{PT}$ with 95\% bootstrap CIs, against the chance floor (grey: random
sets of the observed sizes, bar to the 95th percentile, tick at the mean) and
the P circuit's split-half self-agreement (blue: 2.5th--97.5th percentile, tick
at the median). Hollow: objects kept by both conditions.}
\label{fig:overlap}
\end{figure}

\subsection{Shared selection, separate causes}\label{sec:h2}
Both primary cross-ablation tests are null (Fig.~\ref{fig:h2}). Ablating $\HP$
defeats the typographic override on \hTwoPtoTObs\% of T objects [\hTwoPtoTLo,
\hTwoPtoTHi] (random: \hTwoPtoTNull\%, $p=\hTwoPtoTP$). Ablating $\HT$ defeats the
prior override on \hTwoTtoPObs\% of P objects [\hTwoTtoPLo, \hTwoTtoPHi] (random:
\hTwoTtoPNull\%, $p=\hTwoTtoPP$). Both are \emph{``not distinguishable from
layer-matched random heads at this $n$''}. Each set does defeat its own override:
\hTwoPtoPObs\% and \hTwoTtoTObs\%. \hTwoMeanSentence

The decomposition explains the nulls. The \nShared{} shared heads alone defeat
the prior override on \hTwoStoPObs\% of P objects, but the typographic override
on \hTwoStoTObs\% of T objects (\hTwoStoTCount{} of \hTwoStoTN). The T override
runs through the \nTonly{} T-only heads, which alone defeat it on
\hTwoTOtoTObs\%. The P-only heads defeat it on \hTwoPOtoTObs\% ($p=\hTwoPOtoTP$),
which does not reach $\alpha$. $\HT$ on P reads null for a different reason than
$\HP$ on T. The T-only heads push P's source runs \emph{toward} the image's colour
(mean logit shift \hTwoTOonPShift, against \hTwoStoPShift{} for the shared
heads), so ablating all of $\HT$ cancels the shared heads' effect.

\begin{figure*}[t]
\centering
\includegraphics[width=0.66\textwidth]{fig_cross_ablation}
\caption{Cross-ablation. Share of objects whose override is defeated by
zero-ablating each head set, with 95\% bootstrap CIs. Grey: \hTwoDraws{}
layer-matched random sets (bar to the 95th percentile, tick at the mean). Hollow
markers have no random control (descriptive).}
\label{fig:h2}
\end{figure*}

\subsection{The T-only heads move colour from the image}\label{sec:h3}
On \hThreeNT{} T objects (\hThreeDropT{} dropped because their answers changed
under eager attention or became non-finite), both primary mechanism tests pass
(Fig.~\ref{fig:h3}a,b). The T-only heads place \hThreeAObs{} of their last-token
attention on note tokens [\hThreeALo, \hThreeAHi], against \hThreeANull{} for
random heads ($p=\hThreeAP$). Their direct effect toward the note's colour is
\hThreeWObs{} logits [\hThreeWLo, \hThreeWHi], against \hThreeWNull{}. The
pre-registered mechanism sentence is \emph{``\hThreeMechanism''}. The note makes
up \noteShare\% of the image tokens, but receives \enrichInc{} times its share of
the T-only heads' image attention with a colour word and \enrichNeu{} times with
a neutral word. The heads read text in the image regardless of its content, and
read it more when it names a colour (rise \hThreeRiseObs, random
\hThreeRiseNull).

The shared heads neither attend to the note (\hThreeASObs, random \hThreeASNull,
$p=\hThreeASP$) nor write its colour (\hThreeWSObs{} logits), consistent with
their null ablation on T. The P-only heads attend to the note (\hThreeAPOObs,
random \hThreeAPONull) but do not write its colour (\hThreeWPOObs).

With no note, the T-only heads do the same thing with the object's own colour
(Fig.~\ref{fig:h3}c). On the clean original image they attend to the recoloured
object region (\hThreeObjCleanObs, random \hThreeObjCleanNull) and write the
image's colour (\hThreeColCleanObs{} logits, random \hThreeColCleanNull). On P's
Visual-prompt runs they do likewise (\hThreeObjPVObs; \hThreeColPVObs{} logits).
As a check on the attention measurement, $\HP$ attends to the image more under the
Visual than the Prior prompt, the direction Lietzow et al.\ report (difference
\hThreeRoutingObs, random \hThreeRoutingNull, $p=\hThreeRoutingP$).

\begin{figure*}[t]
\centering
\includegraphics[width=\textwidth]{fig_h3}
\caption{Mechanism. (a) Last-token attention mass on note tokens, mean over the
set's heads; (b) the set's summed direct effect toward the note's colour, both on
T objects with the incongruent note. (c) The T-only heads' direct effect with the
note, and with no note on T's clean images and P's Visual-prompt runs (negative =
toward the image's colour). Grey as in Fig.~\ref{fig:h2}; whiskers are 95\%
bootstrap CIs.}
\label{fig:h3}
\end{figure*}

% ===========================================================================
\section{Discussion}\label{sec:discussion}

\paragraph{Selection is not cause.}
Patching selected the shared heads in both conditions, yet ablating them affects
only P, and under T they neither attend to the note nor write its colour.
Patching asks whether a head's output \emph{carries} the difference between two
runs; ablation asks whether the answer \emph{needs} it. For the shared heads
under T, the two disagree. We do not yet know what patching picks up in them;
redundancy, or an indirect path that ablation does not reveal, are both
possible.\todo{one sentence on what would distinguish these}
Overlap statistics computed from patching alone would have overstated how much
the two overrides share.

\paragraph{Two routes to an override.}
The prior-override heads let prompted knowledge beat the image. The typographic
attack does not use them. It uses heads that, by our measurements, move colour
from wherever it appears in the image to the answer. A note is effective because
it hands those heads a colour, not because it triggers the prior-override
machinery. This fits the attack's dependence on placement: a note on the object
fools more often than one beside it. Yet ablating the T-only heads barely affects ordinary
answers (clean accuracy kept at \hTwoCleanTO\%, random \hTwoCleanRandom\%), so
visual colour must also reach the answer by other routes. This is consistent with
visual grounding being the default.

\paragraph{Limitations.}
One model, quantised to 4 bits. Interventions and attention are measured at the
last token only. P and T share the colour pair but swap its roles, so heads that
\emph{write} an answer may differ partly for that reason. The planned
image-varying contrast (V), whose roles match P, addresses this.\todo{V} Head
selection makes a positive direct effect for the T-only heads partly expected;
the attention result does not share this problem. ``Object'' tokens are the
recoloured pixels, which can omit parts of an object the recolouring left
unchanged. Many tests sit at the smallest p-value 100 random sets allow, so the
excess over random sets and its interval are the effect sizes. A second model
(PaliGemma, whose routing Lietzow et al.\ find to differ) is future work.

\section{Conclusion}
Prompted priors and in-image text both override what a VLM sees, and patching
selects partly the same heads for both. Causally, however, they use different
routes. The prior-override heads are not needed for the typographic attack,
which instead runs through heads that carry colour from the image, whether
painted on the object or written on a note.

\bibliographystyle{plainnat}
\bibliography{refs}

\end{document}
